\documentclass[11pt]{article}
\usepackage[utf8]{inputenc}
\usepackage{minimalist}

\title{\textbf{Simulacija Šumskog Požara}\\[4pt]
       \large Matrični model širenja vatre}
\author{iobih škola — Radionica}
\date{}

\begin{document}
\maketitle
\thispagestyle{fancy}
\markboth{Simulacija Šumskog Požara}{}

Posmatramo šumu kao mrežu ćelija u kojoj se vatra širi korak po korak.
Svaka ćelija ima stanje koje se mijenja prema jednostavnim pravilima.

% ─────────────────────────────────────────────────────────────────────────
\section{Stanja ćelija}

Svaka ćelija u matrici \texttt{karta[r][k]} sadrži jedan karakter:

\begin{center}
\begin{tabular}{@{} c l l @{}}
    \toprule
    \textbf{Znak} & \textbf{Stanje} & \textbf{Opis} \\
    \midrule
    \texttt{.} & Prazno     & Golo tlo — vatra se ne može širiti \\
    \texttt{T} & Drvo       & Živo drvo, može se zapaliti \\
    \texttt{*} & Gori       & Drvo gori ovaj vremenski korak \\
    \texttt{x} & Izgorjelo  & Pepeo, vatra je završena \\
    \bottomrule
\end{tabular}
\end{center}

% ─────────────────────────────────────────────────────────────────────────
\section{Pravila simulacije}

Na svakom vremenskom koraku primjenjuju se dva pravila:

\begin{enumerate}
    \item Svaka ćelija koja \textbf{gori} (\texttt{*}) postaje \textbf{izgorjela} (\texttt{x}).
    \item Svako \textbf{drvo} (\texttt{T}) pored gorućeg susjeda pali se
          sa vjerovatnoćom \texttt{SIRENJE\%}.
\end{enumerate}

\subsection{Ključni dio koda}

\begin{lstlisting}
if (karta[r][k] != GORI) continue;

sljedeca[r][k] = IZGORJELO;

for (int dr = -1; dr <= 1; dr++) {
    for (int dk = -1; dk <= 1; dk++) {
        int nr = r + dr;
        int nk = k + dk;
        if (karta[nr][nk] == DRVO && rand() % 100 < SIRENJE)
            sljedeca[nr][nk] = GORI;
    }
}
\end{lstlisting}

\begin{napomena}
    Važno je ažurirati \texttt{sljedeca}, a čitati iz \texttt{karta}.
    Ako bismo mijenjali \texttt{karta} direktno, nova vatra bi odmah
    uticala na susjedne ćelije u \emph{istom} koraku.
\end{napomena}

% ─────────────────────────────────────────────────────────────────────────
\section{Zadaci}

\begin{zadatak}{1}
Postavi \texttt{GUSTOCA = 30}. Kako se ponaša vatra u rijetkoj šumi?
Postoji li gustoća pri kojoj se vatra ugasi sama od sebe?
\end{zadatak}

\begin{zadatak}{2}
Šta se desi ako postaviš \texttt{SIRENJE = 100}? A \texttt{SIRENJE = 0}?
\end{zadatak}

\begin{zadatak}{3}
Dodaj horizontalni red \emph{praznih} ćelija koji prolazi cijelom širinom
mape (npr.\ red 10). Može li vatra preskočiti taj red?
\end{zadatak}

\end{document}
